
\paragraph{\textbf{Theorem 1.5.}}
\SGARefTarget{sga3:xx:1.5}
\emph{Let \((G,T,\alpha)\) be an \(S\)-elementary system.}
\begin{enumerate}
\item[\textup{\SGAInlineRefTarget{sga3:xx:1.5:item:i}(i)}] \emph{There exists a unique morphism of groups
  with operator group \(T\)
  \[
    \exp\colon W(\mathfrak g^\alpha)\longrightarrow G
  \]
  whose induced morphism on Lie algebras is the canonical map
  \(\mathfrak g^\alpha\to\mathfrak g\).\footnote{It will be seen
  below, in \SGARefLink{sga3:xx:5.9}{Corollary 5.9}, that \(\exp\) is an isomorphism from
  \(W(\mathfrak g^\alpha)\) onto a closed subgroup of \(G\).}
  Equivalently, \(\exp\) is the unique morphism such that, for every
  \(S'\to S\), \(t\in T(S')\), and
  \(X,X'\in W(\mathfrak g^\alpha)(S')\),
  \[
  \begin{aligned}
    \exp(X+X')&=\exp(X)\exp(X'),\\
    \intaut(t)\bigl(\exp(X)\bigr)&=\exp\bigl(\alpha(t)X\bigr),\\
    \Lie(\exp)(X)&=X.
  \end{aligned}
  \]}
\item[\textup{\SGAInlineRefTarget{sga3:xx:1.5:item:ii}(ii)}] \emph{Define analogously, in the
  \(S\)-elementary system \((G,T,-\alpha)\),
  \[
    \exp\colon W(\mathfrak g^{-\alpha})\longrightarrow G.
  \]
  Then the morphism
  \[
    W(\mathfrak g^{-\alpha})\times_S T
      \times_S W(\mathfrak g^\alpha)\longrightarrow G,
    \qquad
    (Y,t,X)\longmapsto \exp(Y)\,t\,\exp(X),
  \]
  defined on points by the displayed formula, is an open
  immersion.}
\end{enumerate}

Suppose for the moment that the required exponential morphisms
exist, and prove the other assertions. We first prove \textup{\SGARefLink{sga3:xx:1.5:item:ii}{(ii)}}.
Both sides are of finite presentation and flat over \(S\), so it is
enough to work when \(S\) is the spectrum of an algebraically closed
field \textup{(SGA 1, I 5.7 and VIII 5.5)}. Put \(S=\Spec(k)\), and
choose
\[
  X\in\Gamma(S,\mathfrak g^\alpha)^\times,
  \qquad
  Y\in\Gamma(S,\mathfrak g^{-\alpha})^\times.
\]
It is enough to prove that
\[
  \mathbb G_{a,k}\times_k T\times_k\mathbb G_{a,k}
    \longrightarrow G,
  \qquad
  (y,t,x)\longmapsto \exp(yY)\,t\,\exp(xX),
\]
is an open immersion. By \SGARefLink{sga3:xx:section:1}{1.1} and \SGARefLink{sga3:xx:1.2}{1.2}, there exist \(a,b\in k\)
such that
\[
  \exp(yY)=p_{-\alpha}(ay),
  \qquad
  \exp(xX)=p_\alpha(bx).
\]
Since
\(\exp\colon W(\mathfrak g^{-\alpha})\to G\) induces a monomorphism
on Lie algebras, \(a\neq0\); likewise \(b\neq0\). The assertion is
therefore reduced to \SGARefLink{sga3:xx:section:1}{1.1}.

The uniqueness of \(\exp\) may be proved locally on \(S\). We reduce
to the case in which \(\mathfrak g^\alpha\) and
\(\mathfrak g^{-\alpha}\) are free, and then apply \SGARefLink{sga3:xx:1.2}{1.2} with
\(Q=T\) and \(q=1\).

It remains to prove existence. Faithfully flat descent, together
with the uniqueness just proved, shows that it is enough to prove
existence locally on \(S\) for the fpqc topology. The usual
finite-presentation arguments reduce first to \(S\) noetherian and
then to \(S\) noetherian local. We may therefore assume that
\[
  S=\Spec(A),
\]
where \(A\) is a complete noetherian local ring with algebraically
closed residue field \(k\). Let
\[
  p_0\colon\mathbb G_{a,k}\longrightarrow G_k
\]
be a monomorphism of \(k\)-groups normalized by \(T_k\), with
multiplier
\(\alpha_0=\alpha\otimes_A k\); such a morphism exists by \SGARefLink{sga3:xx:section:1}{1.1}.
By \SGARefLink{sga3:xx:section:1}{1.1} and \SGARefLink{sga3:xx:1.2}{1.2}, the associated morphism
\[
  T_k\mathbin{\cdot_{\alpha_0}}\mathbb G_{a,k}
    \longrightarrow G_k
\]
is an immersion and hence, in particular, a monomorphism. We
provisionally admit the following two lemmas.

\paragraph{\textbf{Lemma 1.6.}}
\SGARefTarget{sga3:xx:1.6}
\emph{Let \(S\) be a scheme, \(G\) an \(S\)-group of finite
presentation, \(T\) an \(S\)-torus, \(\alpha\) a character of \(T\)
that is nontrivial on every fiber, and \(s_0\in S\). Let
\[
  f\colon T\mathbin{\cdot_\alpha}\mathbb G_{a,S}\longrightarrow G
\]
be an \(S\)-group morphism such that \(f_{s_0}\) is a monomorphism
and the restriction of \(f\) to \(T\) is a monomorphism. There
exists an open neighborhood \(U\) of \(s_0\) such that
\(f|_U\) is a monomorphism.}

\paragraph{\textbf{Lemma 1.7.}}
\SGARefTarget{sga3:xx:1.7}
\emph{Let \(A\) be a complete noetherian local ring with
algebraically closed residue field \(k\), let
\((G,T,\alpha)\) be an \(A\)-elementary system, and let
\[
  p_0\colon\mathbb G_{a,k}\longrightarrow G_k
\]
be a morphism of \(k\)-groups normalized by \(T_k\), with
multiplier \(\alpha\otimes_A k\). There exists a group morphism
\[
  p\colon\mathbb G_{a,A}\longrightarrow G
\]
normalized by \(T\), with multiplier \(\alpha\).}

Let \(p\) be the morphism asserted by \SGARefLink{sga3:xx:1.7}{1.7}, and let
\[
  f\colon T\mathbin{\cdot_\alpha}\mathbb G_{a,S}\longrightarrow G
\]
be the corresponding morphism. It satisfies the hypotheses of \SGARefLink{sga3:xx:1.6}{1.6},
so it is a monomorphism; in particular, \(p\) is a monomorphism.
The theorem now follows from \SGARefLink{sga3:xix:corollary:4:9}{Exposé XIX, 4.9}.

\medskip
\noindent\emph{Proof of \SGARefLink{sga3:xx:1.6}{1.6}.}
Let
\[
  \epsilon\colon S\longrightarrow
    T\mathbin{\cdot_\alpha}\mathbb G_{a,S}
\]
be the identity section. Since \(f\) is unramified at
\(\epsilon(s_0)\), it is unramified at \(\epsilon(s)\) for all
\(s\) in an open neighborhood \(U\) of \(s_0\). Hence \(f|_U\) is
unramified \textup{(\SGARefLink{sga3:x:lemma:3:5}{Exposé X, 3.5})},\footnote{See also Exposé
VI\(_B\), \SGARefLink{sga3:xx:2.5}{2.5}.} and its kernel \(\Ker(f)_U\) is unramified over
\(U\). To prove that \(f|_U\) is a monomorphism, it is therefore
enough\footnote{By EGA IV\(_4\), 17.9.1.} to prove that
\(\Ker(f)_U\) is radicial over \(U\), which is a set-theoretic
question. We are reduced to the following assertion.

\paragraph{\textbf{Lemma 1.8.}}
\SGARefTarget{sga3:xx:1.8}
\emph{Let \(k\) be an algebraically closed field, and let \(N\) be
an invariant subgroup of
\(T\mathbin{\cdot_\alpha}\mathbb G_{a,k}\), where \(\alpha\) is a
nontrivial character of the torus \(T\). Suppose that \(N\) is
étale over \(k\) and that \(N\cap T=\{e\}\). Then \(N=\{e\}\).}

We have
\[
  \intaut(t')(t,x)=\bigl(t,\alpha(t')x\bigr).
\]
If \((t,x)\) is a point of \(N\) with \(x\neq0\), then
\((t,zx)\) is also a point of \(N\) for every \(z\in k^\times\).
Thus \((t,x)\) is not isolated, contradicting the quasi-finiteness
of \(N\). Set-theoretically, therefore, \(N\subset T\), and the
hypothesis \(N\cap T=\{e\}\) completes the proof.
